\begin{lemma}[Quantitative spanning super-regular thinning]
Let $0<\beta<\eta$, $0<\rho<D\leq1$, and $\xi>0$. Put $\lambda=\rho/D$ and assume
\[
4(\lambda\beta+\xi)\leq\eta\quad\text{and}\quad 10\eta<\rho.
\]
There is $n_0$ such that the following holds. If $L,R$ are disjoint, both have size at least $n_0$, their sizes differ by at most one, and $G[L,R]$ is $(\beta,D)$-super-regular in the sense of \noderef{SuperRegularPair}, then $G$ contains a spanning subgraph $G'$ for which $G'[L,R]$ is $(\eta,\rho)$-super-regular.
\end{lemma}
\begin{proof}
First apply \noderef{ThinningRealParameters}: all parameters are finite. Indeed $D\leq1$, $\rho<D$,
$\eta\leq10\eta<\rho$, and $\beta<\eta$ establish finiteness of
$D,\rho,\eta,\beta$. Also
$\xi\leq\lambda\beta+\xi\leq4(\lambda\beta+\xi)\leq\eta$ establishes
finiteness of $\xi$. Positivity and finiteness permit real conversions
throughout the following argument. Put $q=\rho/D$ in the real numbers;
then $0<q<1$, $qD=\rho$, and the budget hypothesis gives
$q\beta+\xi\leq\eta/4\leq\eta$. Also $\beta<\eta<\rho<D$.

Apply \noderef{NearEqualRetentionThreshold} with the positive real
conversions of $\xi,\eta$ to choose $n_0\geq1$, before choosing any
vertex type or graph. On a graph with admissible parts, use
\noderef{IndependentEdgeSample} on $V\times V$ with probability $q$.
This is one probability measure by \noderef{IndependentUnequalEdgeSampleLaw}
specialized to the constant coordinate probabilities $q_i=q$.
By \noderef{RetentionBadEventUnion}, the event failing
\noderef{RetentionCountControl} has measure at most
\[
\operatorname{ofReal}\bigl(
2(|L|+|R|)e^{-2\xi^2\min(|L|,|R|)}
+2\,2^{|L|+|R|}e^{-2\xi^2\eta^2|L||R|}\bigr)<1.
\]
The strict inequality is the chosen numerical threshold and uses both
near-equal inequalities. If no outcome had count control, this failure
event would be the whole sample space and have measure one, a
contradiction. Choose an outcome $\omega$ with count control, and set
$G'=G_\omega$ using \noderef{CrossEdgeRetention}. This graph has the
original vertex set. By \noderef{CrossEdgeRetentionCounts} it is
contained in $G$, and its rectangle and degree counts are exactly
the retained coordinate counts tested in count control.

For $v\in L$ take $m=|R|$, $a=\deg_G(v;R)$ and
$b=\deg_{G'}(v;R)$. The original super-regularity hypothesis supplies
$(D-\beta)m\leq a\leq(D+\beta)m$ by \noderef{SuperRegularPair}.
The coordinate bijection $w\mapsto(v,w)$ shows that the original
coordinate set has size $a$, and the degree identity just cited gives
$|b-qa|\leq\xi m$ from count control. Apply
\noderef{ThinningCountInterval} to obtain
$(\rho-\eta)|R|\leq b\leq(\rho+\eta)|R|$.
For $w\in R$, use $m=|L|$ and the bijection $v\mapsto(v,w)$;
symmetry identifies the original adjacency count with
$\deg_G(w;L)$. The same transfer proves the right-degree interval.

Finally take $U\subseteq L,W\subseteq R$ with
$|U|>\eta|L|$, $|W|>\eta|R|$. These subsets are nonempty, since
$\eta>0$ and the parts have size at least one. Because $\beta<\eta$,
they also satisfy the original $\beta$-thresholds. Thus the regularity
clause of \noderef{SuperRegularPair} gives
$D-\beta\leq d_G(U,W)\leq D+\beta$.
By \noderef{PairDensity}, $d_G(U,W)=e_G(U,W)/m$, where
$m=|U||W|$ is positive and finite. Multiplication by $m$ preserves
and reflects each inequality (also for extended nonnegative endpoints),
so
\[
(D-\beta)|U||W|\leq e_G(U,W)\leq(D+\beta)|U||W|.
\]
Count control and the exact rectangle identity give
$|e_{G'}(U,W)-qe_G(U,W)|\leq\xi|U||W|$. Apply
\noderef{ThinningCountInterval} with $m=|U||W|$ to obtain
$(\rho-\eta)m\leq e_{G'}(U,W)\leq(\rho+\eta)m$.
Dividing both inequalities by the same positive finite $m$, and using
the density formula of \noderef{PairDensity} for $G'$, gives
$\rho-\eta\leq d_{G'}(U,W)\leq\rho+\eta$.
This is precisely the regularity clause; together with the two degree
intervals it proves \noderef{SuperRegularPair}. Every normalization
used the actual opposite-part size. No equality of the part sizes or
graph-dependent threshold was used.
\end{proof}
